\documentclass[11pt]{article}
\usepackage[margin=0.95in]{geometry}
\usepackage[T1]{fontenc}
\usepackage[utf8]{inputenc}
\usepackage{amsmath,amssymb,amsthm}
\usepackage{xcolor}
\usepackage{microtype}
\usepackage[colorlinks=true,linkcolor=blue!55!black,citecolor=blue!55!black,urlcolor=blue!55!black]{hyperref}

\newtheorem{theorem}{Theorem}
\newtheorem{lemma}{Lemma}
\newtheorem{remark}{Remark}
\newcommand{\Hh}{\mathcal H}
\newcommand{\N}{\mathbb N}
\newcommand{\T}{T}

\title{Strong Gibbs convergence for every $\ell^2$ initial condition}
\author{Candidate full answer to Remark (2) in arXiv:2011.00457}
\date{12 August 2026}

\begin{document}
\maketitle

\begin{center}
\fbox{\parbox{0.92\textwidth}{\raggedright\textbf{Status: candidate full answer.}
Under the hypotheses of the source's Main Theorem, the untruncated master-equation
semigroup converges strongly to its Gibbs-mode projection for every initial
condition in $\ell^2$.  A quantitative head--tail estimate is available, but
operator-norm convergence is impossible at every positive time.}}
\end{center}

\section{Source question and answer}

B\"ogli and Vuillermot~\cite{source} study a compact, nonnormal generator $A$
on $\ell^2_{\mathbb C}$.  Under the gap-control hypotheses of their Main
Theorem, its eigenvectors $(\widehat p_k)_{k\geq1}$ form a Schauder basis and
\begin{equation}\label{eq:spectral}
 e^{tA}x=\sum_{k=1}^{\infty}
 \langle x,\widehat q_k\rangle e^{t\nu_k}\widehat p_k
 \qquad(t\geq0)
\end{equation}
in norm.  Here $\nu_1=0$, while
\begin{equation}\label{eq:order}
 \nu_2<\nu_3<\cdots<0,\qquad \nu_k\longrightarrow0,
\end{equation}
and $\widehat p_1$ is a scalar multiple of the Gibbs vector.

Remark (2) after their corollary proves the limit below for $x\in\ell^1$ and
asks whether it holds for every $x\in\ell^2$.  Define the bounded coordinate
projection
\[
 \Pi x=\langle x,\widehat q_1\rangle\widehat p_1.
\]

\begin{theorem}[Full $\ell^2$ convergence]\label{thm:main}
Under the hypotheses of the source's Main Theorem,
\begin{equation}\label{eq:limit}
 \lim_{t\to\infty}\|e^{tA}x-\Pi x\|_2=0
 \qquad\text{for every }x\in\ell^2_{\mathbb C}.
\end{equation}
More precisely, let
\[
 P_Nx=\sum_{k=1}^N\langle x,\widehat q_k\rangle\widehat p_k,
 \qquad K=\sup_N\|P_N\|<\infty.
\]
For every $N\geq2$ and $t\geq0$,
\begin{equation}\label{eq:tail}
 \|(e^{tA}-\Pi)x\|_2
 \leq 2K e^{t\nu_N}\|x\|_2
 +(1+K)\|(I-P_N)x\|_2.
\end{equation}
On the other hand,
\begin{equation}\label{eq:sharp}
 \|e^{tA}-\Pi\|\geq1\qquad(t\geq0),
\end{equation}
so the convergence is never in operator norm and admits no uniform decay rate
on the unit ball.
\end{theorem}

\section{Proof intuition}

Coefficientwise convergence in \eqref{eq:spectral} is not enough by itself:
the eigenvectors are a Schauder basis, not asserted to be unconditional.
The missing control comes from the order of the eigenvalues.  For fixed $t$,
the non-Gibbs multipliers
\[
 e^{t\nu_2}\leq e^{t\nu_3}\leq\cdots\longrightarrow1
\]
are monotone, so their total variation is uniformly bounded.  Abel summation
transfers that scalar variation bound to the basis multiplier through the
uniformly bounded partial-sum projections $P_N$.  The resulting operators are
uniformly bounded in time.  They vanish on each finite basis span as
$t\to\infty$, and density finishes the proof.

The same calculation separates a finite head, which decays at rate
$e^{t\nu_N}$, from an arbitrarily small basis tail.  Yet high-index
eigenvalues are arbitrarily close to zero, and their multipliers are
arbitrarily close to one at each fixed time; this gives the sharp norm
obstruction \eqref{eq:sharp}.

\section{Abel multiplier lemma}

\begin{lemma}\label{lem:abel}
Let $(u_k)$ be a Schauder basis of a Banach space, with partial-sum
projections $P_N$ and $K=\sup_N\|P_N\|$.  If $(a_k)$ has finite total
variation and $a_k\to L$, then its basis multiplier satisfies
\begin{equation}\label{eq:abel}
 M_a=L I+\sum_{n=1}^{\infty}(a_n-a_{n+1})P_n,
 \qquad
 \|M_a\|\leq |L|+K\sum_{n=1}^{\infty}|a_n-a_{n+1}|.
\end{equation}
\end{lemma}

\begin{proof}
For every finite basis expansion, the coefficient of $u_k$ on the right of
\eqref{eq:abel} is
\[
 L+\sum_{n=k}^{\infty}(a_n-a_{n+1})=a_k.
\]
The series of operators converges absolutely by the variation hypothesis,
and the norm estimate follows from $\|P_n\|\leq K$.  Density extends the
identity to the whole space.
\end{proof}

\section{Proof of Theorem~\ref{thm:main}}

Put $D_t=e^{tA}-\Pi$.  Relative to the basis $(\widehat p_k)$, its scalar
multipliers are
\[
 a_1(t)=0,\qquad a_k(t)=e^{t\nu_k}\quad(k\geq2).
\]
By \eqref{eq:order}, $a_k(t)$ increases for $k\geq2$ and tends to $1$.
Consequently
\[
 \sum_{n=1}^{\infty}|a_n(t)-a_{n+1}(t)|
 =a_2(t)+(1-a_2(t))=1.
\]
Lemma~\ref{lem:abel} gives the time-uniform bound
\begin{equation}\label{eq:uniform}
 \|D_t\|\leq1+K.
\end{equation}

To obtain the stated estimate, truncate the multiplier after index $N$.
The finite scalar sequence
\[
 0,a_2(t),\ldots,a_N(t),0,0,\ldots
\]
has total variation $2a_N(t)=2e^{t\nu_N}$ and limit zero.  A second application
of Lemma~\ref{lem:abel}, followed by \eqref{eq:uniform}, yields
\begin{align*}
 \|D_tx\|_2
 &\leq\|D_tP_Nx\|_2+\|D_t(I-P_N)x\|_2\\
 &\leq2K e^{t\nu_N}\|x\|_2+(1+K)\|(I-P_N)x\|_2,
\end{align*}
which is \eqref{eq:tail}.

Given $x$ and $\varepsilon>0$, first choose $N$ so large that the tail term
in \eqref{eq:tail} is below $\varepsilon/2$.  Since $\nu_N<0$, next choose
$t$ so large that the head term is below $\varepsilon/2$.  This proves
\eqref{eq:limit}.

Finally, $D_t\widehat p_k=e^{t\nu_k}\widehat p_k$ for $k\geq2$.  Hence
\[
 \|D_t\|\geq\sup_{k\geq2}e^{t\nu_k}=1
\]
because $\nu_k\to0$.  This proves \eqref{eq:sharp} and completes the proof.

\section{Scope and literature check}

The theorem answers exactly the ordinary-$\ell^2$ question under all the
assumptions used by the source to obtain its Schauder-basis spectral theorem.
It does not remove those hypotheses.  The estimate \eqref{eq:tail} gives an
orbit-dependent convergence mechanism, while \eqref{eq:sharp} explains why
no spectral-gap-style uniform exponential estimate can exist.

Searches through 12 August 2026 located the 2022 journal version of the
source and a later grand-canonical paper by the same authors in a different,
weighted self-adjoint sequence space, but no statement of this exact
nonnormal ordinary-$\ell^2$ conclusion.  Novelty should nevertheless be
checked in specialist citation databases before publication.

\begin{thebibliography}{9}
\bibitem{source}
S. B\"ogli and P.-A. Vuillermot,
\emph{A spectral theorem for the semigroup generated by a class of infinitely
many master equations}, Acta Appl. Math. 178 (2022), article 4;
arXiv:2011.00457.
\end{thebibliography}

\end{document}
